\begingroup
\clearpage
\let\clearpage\relax
\vspace*{-16pt}
\chapter[PROBLEM STATEMENT]{PROBLEM STATEMENT}\label{sec.probstmnt}
\endgroup

\comm{BB is here.}

This chapter formalizes the optimization problem studied in this
dissertation and establishes the notation used throughout. EHR data
are assumed to supply all inputs needed for the optimization model.
Let $P$ be the set of patients in the EHR dataset, and let
$\widetilde{P} \subseteq P$ be the subset of patients recorded as
expired. The set of diseases (or disease categories) $V$ consists of
standardized codes drawn from the International Classification of
Diseases (ICD)~\citep{cdc_icd9cm_2013, cdc_icd10cm_2022}. For each
disease $u \in V$, let $A_u \subseteq P$ be the set of patients
diagnosed with $u$, and for each patient $i \in P$, let
$D_i \subseteq V$ be the set of diseases with which patient $i$ has
been diagnosed. Without loss of generality, we assume that $A_u$ and
$D_i$ are nonempty for all $u \in V$ and $i \in P$.

From the same EHR dataset, we construct a \textit{comorbidity graph}
whose vertices are diseases and whose edges capture pairwise
co-occurrence among diseases in the patient population $P$. Following
prior work on comorbidity
graphs~\citep{KALGOTRA2020,VaghfiMohebbi2025lethalclique}, we use the
Ochiai coefficient (akin to cosine similarity for binary vectors) to
quantify co-occurrence as
\begin{equation}
    \label{eq.SCI}
    \sigma(u,v) = \frac{|A_u \cap A_v|}{\sqrt{|A_u| \times |A_v|}},
\end{equation}
which normalizes the co-occurrence of a pair of diseases by the
geometric mean of their individual occurrences in
$P$~\citep{ochiai1957zoogeographic,Salton}. For a
user-specified threshold $\Delta \in [0,1]$, the edge set of the
comorbidity graph $G=(V,E)$ is
\begin{equation}
    E \coloneqq \left\{ \{u,v\} \in \binom{V}{2} :
        \sigma(u,v) \ge \Delta \right\},
\end{equation}
where $\binom{V}{2}$ denotes the set of all two-element subsets of
$V$. A clique is a subset of pairwise adjacent vertices; in our
setting, cliques model tightly knit groups of diseases that co-occur
frequently across the patient population.

\begin{definition}\label{def.MR}
Given a clique of diseases $C \subseteq V$, its \textit{mortality
rate} $\mu(C)$ is defined as
\begin{equation}
    \label{eq.MR}
    \mu(C) \coloneqq \frac{\left|\num(C)\right|}{\left|\den(C)\right|},
\end{equation}
where $\den(C) \coloneqq \bigcap_{u \in C} A_u$ and
$\num(C) \coloneqq \den(C) \cap \widetilde{P}$.
\end{definition}

The denominator $|\den(C)|$ counts the patients in $P$ diagnosed with
every disease in $C$, and the numerator $|\num(C)|$ counts those
patients who have expired. We set $\mu(C) = 0$ whenever $C$ or
$\den(C)$ is empty.

\begin{definition}\label{def.lfreq}
For a positive integer $\ell$, a clique $C$ is
\textit{$\ell$-frequent} if $|\den(C)| \ge \ell$, and
\textit{non-$\ell$-frequent} otherwise. A non-$\ell$-frequent clique
$C$ is called \textit{minimal} if $|\den(C)| \le \ell - 1$ and
$|\den(C \setminus \{u\})| \ge \ell$ for every $u \in C$.
\end{definition}

The \textit{maximum mortality rate clique problem} (MMRCP) seeks
\begin{equation}
    \max\bigl\{\mu(C) : C \text{ is a nonempty $\ell$-frequent
        clique in } G\bigr\}.
    \label{eq.MMRCP}
\end{equation}

The precise meaning of being recorded as expired depends on how the EHR dataset defines $\widetilde{P}$. For instance, a patient may be considered
expired if they died or were discharged to hospice care within a given
time window after their last clinical encounter.

Throughout this dissertation, we apply a simple preprocessing step to
reduce the size of the comorbidity graph before solving MMRCP. Since
only $\ell$-frequent cliques are of interest, any vertex $u$ with
$|A_u| \le \ell - 1$ can be removed from $G$, as it cannot belong to
any $\ell$-frequent clique. Likewise, any edge $\{u,v\}$ with
$|A_u \cap A_v| \le \ell - 1$ can be removed, since no clique
containing both $u$ and $v$ can be $\ell$-frequent.

The monotonicity properties of MMRCP and its computational complexity
are discussed in Chapter~\ref{sec.theory}.





